% Auto-generated from raw.txt
% Usage:
%   \vocabentry{word}{/ipa/ (pos)}{definition}{example}{note}

\vocabentry{Administration}
{/ədˌmɪnɪˈstreɪʃn/ (n)}
{The process or activity of running a business, organization, or government.}
{The current administration is focusing on economic recovery and healthcare reform.}
{Đây là một danh từ dùng để chỉ bộ máy hành chính hoặc chính quyền của một nhiệm kỳ cụ thể.}

\vocabentry{Ambassador}
{/æmˈbæsədər/ (n)}
{An accredited diplomat sent by a country as its official representative to a foreign country.}
{The American ambassador met with the foreign minister to discuss trade agreements.}
{Đây là một danh từ dùng để chỉ đại sứ, người đại diện cao nhất của một quốc gia tại nước ngoài.}

\vocabentry{Amendment}
{/əˈmendmənt/ (n)}
{A minor change or addition designed to improve a text, piece of legislation, etc.}
{The First Amendment to the Constitution protects the freedom of speech and religion.}
{Đây là một danh từ dùng để chỉ sự sửa đổi hoặc bổ sung cho một đạo luật hoặc hiến pháp.}

\vocabentry{Anarchy}
{/ˈænərki/ (n)}
{A state of disorder due to absence or nonrecognition of authority.}
{Without a stable government, the country could quickly descend into anarchy.}
{Đây là một danh từ dùng để chỉ tình trạng vô chính phủ hoặc hỗn loạn.}

\vocabentry{Aristocracy}
{/ˌærɪˈstɑːkrəsi/ (n)}
{The highest class in certain societies, especially those holding hereditary titles or offices.}
{Members of the aristocracy traditionally held a significant amount of political power.}
{Đây là một danh từ dùng để chỉ tầng lớp quý tộc hoặc chính thể quý tộc.}

\vocabentry{Authority}
{/əˈθɔːrəti/ (n)}
{The power or right to give orders, make decisions, and enforce obedience.}
{The local housing authority is responsible for managing public apartment buildings.}
{Đây là một danh từ dùng để chỉ quyền lực, uy quyền hoặc cơ quan có thẩm quyền.}

\vocabentry{Autocracy}
{/ɔːˈtɑːkrəsi/ (n)}
{A system of government by one person with absolute power.}
{The nation moved away from autocracy towards a more democratic system of government.}
{Đây là một danh từ dùng để chỉ chính thể độc tài, chuyên chế.}

\vocabentry{Ballot}
{/ˈbælət/ (n)}
{A process of voting, in writing and typically in secret.}
{The voters cast their ballots to elect the new members of parliament.}
{Đây là một danh từ dùng để chỉ phiếu bầu hoặc quá trình bỏ phiếu.}

\vocabentry{Bill}
{/bɪl/ (n)}
{A draft of a proposed law presented to parliament for discussion.}
{The government introduced a new bill to increase funding for public schools.}
{Đây là một danh từ dùng để chỉ dự thảo luật trước khi được thông qua thành luật chính thức.}

\vocabentry{Bipartisan}
{/ˌbaɪˈpɑːrtɪzn/ (adj)}
{Involving the agreement or cooperation of two political parties that usually oppose each other's policies.}
{There is bipartisan support for the new legislation on environmental protection.}
{Đây là một tính từ dùng để chỉ sự đồng thuận của cả hai đảng phái chính trị.}

\vocabentry{Bureaucracy}
{/bjʊˈrɑːkrəsi/ (n)}
{A system of government in which most of the important decisions are made by state officials rather than by elected representatives.}
{Many people complain about the slow and complicated bureaucracy of the government.}
{Đây là một danh từ dùng để chỉ bộ máy quan liêu hoặc hệ thống hành chính.}

\vocabentry{Cabinet}
{/ˈkæbɪnət/ (n)}
{A body of advisers to the head of a state who also serve as the heads of government departments.}
{The President met with his Cabinet to discuss national security issues.}
{Đây là một danh từ dùng để chỉ nội các, gồm các bộ trưởng cấp cao.}

\vocabentry{Campaign}
{/kæmˈpeɪn/ (n)}
{An organized course of action to achieve a particular goal, especially a political one.}
{The candidate spent millions of dollars on his television advertising campaign.}
{Đây là một danh từ dùng để chỉ chiến dịch vận động, đặc biệt là chiến dịch tranh cử.}

\vocabentry{Candidate}
{/ˈkændɪdət/ (n)}
{A person who applies for a job or is nominated for election.}
{There are several strong candidates running for the position of mayor.}
{Đây là một danh từ dùng để chỉ ứng cử viên.}

\vocabentry{Coalition}
{/ˌkoʊəˈlɪʃn/ (n)}
{An alliance for combined action, especially a temporary alliance of political parties forming a government.}
{The two smaller parties formed a coalition to gain a majority in parliament.}
{Đây là một danh từ dùng để chỉ sự liên minh giữa các đảng phái.}

\vocabentry{Constitution}
{/ˌkɑːnstɪˈtuːʃn/ (n)}
{A body of fundamental principles or established precedents according to which a state is governed.}
{The nation's constitution outlines the rights and responsibilities of its citizens.}
{Đây là một danh từ dùng để chỉ hiến pháp.}

\vocabentry{Constituency}
{/kənˈstɪtʃuənsi/ (n)}
{A body of voters in a specified area who elect a representative to a legislative body.}
{The MP is working hard to address the concerns of his local constituency.}
{Đây là một danh từ dùng để chỉ khu vực cử tri hoặc đơn vị bầu cử.}

\vocabentry{Corruption}
{/kəˈrʌpʃn/ (n)}
{Dishonest or fraudulent conduct by those in power, typically involving bribery.}
{The government has launched a campaign to eliminate corruption in the civil service.}
{Đây là một danh từ dùng để chỉ sự tham nhũng.}

\vocabentry{Democracy}
{/dɪˈmɑːkrəsi/ (n)}
{A system of government by the whole population or all the eligible members of a state, typically through elected representatives.}
{Freedom of speech is a fundamental principle of a healthy democracy.}
{Đây là một danh từ dùng để chỉ chế độ dân chủ.}

\vocabentry{Dictator}
{/ˈdɪkteɪtər/ (n)}
{A ruler with total power over a country, typically one who has obtained control by force.}
{The dictator maintained control through fear and the suppression of dissent.}
{Đây là một danh từ dùng để chỉ kẻ độc tài.}

\vocabentry{Diplomacy}
{/dɪˈploʊməsi/ (n)}
{The profession, activity, or skill of managing international relations.}
{Diplomacy is essential for resolving conflicts between nations peacefully.}
{Đây là một danh từ dùng để chỉ ngành ngoại giao hoặc sự khôn khéo trong đàm phán.}

\vocabentry{Electorate}
{/ɪˈlektərət/ (n)}
{All the people in a country or area who are entitled to vote in an election.}
{The candidate is trying to appeal to the younger members of the electorate.}
{Đây là một danh từ dùng để chỉ toàn thể cử tri.}

\vocabentry{Executive}
{/ɪɡˈzekjətɪv/ (adj)}
{Having the power to put plans, actions, or laws into effect.}
{The executive branch of the government is responsible for enforcing the laws.}
{Đây là một tính từ hoặc danh từ dùng để chỉ nhánh hành pháp của chính phủ.}

\vocabentry{Federal}
{/ˈfedərəl/ (adj)}
{Relating to a system of government in which several states form a unity but remain independent in internal affairs.}
{The federal government has authority over issues like national defense and foreign policy.}
{Đây là một tính từ dùng để chỉ chính quyền liên bang.}

\vocabentry{Governance}
{/ˈɡʌvərnəns/ (n)}
{The action or manner of governing.}
{Good governance is essential for economic development and social stability.}
{Đây là một danh từ dùng để chỉ sự quản trị hoặc cách thức điều hành chính phủ/tổ chức.}

\vocabentry{Hierarchy}
{/ˈhaɪərɑːrki/ (n)}
{A system or organization in which people or groups are ranked one above the other according to status or authority.}
{There is a strict hierarchy within the military and the government bureaucracy.}
{Đây là một danh từ dùng để chỉ hệ thống cấp bậc trong bộ máy chính quyền.}

\vocabentry{Ideology}
{/ˌaɪdiˈɑːlə dʒi/ (n)}
{A system of ideas and ideals, especially one that forms the basis of economic or political theory and policy.}
{The two political parties have very different ideologies regarding taxation and social welfare.}
{Đây là một danh từ dùng để chỉ hệ tư tưởng chính trị.}

\vocabentry{Inauguration}
{/ɪˌnɔːɡjəˈreɪʃn/ (n)}
{The formal admission of someone to office.}
{Thousands of people gathered in the capital to witness the President's inauguration.}
{Đây là một danh từ dùng để chỉ lễ nhậm chức.}

\vocabentry{Incumbent}
{/ɪnˈkʌmbənt/ (n)}
{The holder of an office or post.}
{The incumbent mayor is facing a tough challenge in the upcoming election.}
{Đây là một danh từ dùng để chỉ người đang đương nhiệm một chức vụ chính trị.}

\vocabentry{Judiciary}
{/dʒuˈdɪʃieri/ (n)}
{The judicial authorities of a country; judges collectively.}
{An independent judiciary is essential for protecting the rule of law.}
{Đây là một danh từ dùng để chỉ hệ thống tư pháp hoặc bộ máy tòa án.}

\vocabentry{Legislation}
{/ˌledʒɪsˈleɪʃn/ (n)}
{Laws, considered collectively.}
{The government is planning to introduce new legislation to tackle climate change.}
{Đây là một danh từ dùng để chỉ hệ thống luật pháp hoặc việc ban hành luật.}

\vocabentry{Legislature}
{/ˈledʒɪsleɪtʃər/ (n)}
{The legislative body of a country or state. (Đây là một danh từ dùng để chỉ cơ quan lập pháp (Quốc hội).).}
{The legislature is responsible for debating and passing new laws.}
{}

\vocabentry{Liberalism}
{/ˈlɪbrəlɪzəm/ (n)}
{A political and social philosophy that promotes individual rights, civil liberties, democracy, and free enterprise.}
{Liberalism emphasizes the importance of personal freedom and government accountability.}
{Đây là một danh từ dùng để chỉ chủ nghĩa tự do.}

\vocabentry{Lobbying}
{/ˈlɑːbiɪŋ/ (n)}
{Seeking to influence a politician or public official on an issue.}
{Many corporations spend large sums of money on lobbying to influence government policy.}
{Đây là một danh từ dùng để chỉ hoạt động vận động hành lang.}

\vocabentry{Monarchy}
{/ˈmɑːnərki/ (n)}
{A form of government with a monarch (king or queen) at the head.}
{Britain is a constitutional monarchy where the Queen has limited political power.}
{Đây là một danh từ dùng để chỉ chế độ quân chủ.}

\vocabentry{Opposition}
{/ˌɑːpəˈzɪʃn/ (n)}
{Resistance or dissent, expressed in action or argument.}
{The leader of the opposition criticized the government's handling of the economy.}
{Đây là một danh từ dùng để chỉ đảng đối lập trong chính trường.}

\vocabentry{Parliament}
{/ˈpɑːrləmənt/ (n)}
{The highest legislature, consisting of the Sovereign, the House of Lords, and the House of Commons.}
{The law was passed by parliament after a lengthy and heated debate.}
{Đây là một danh từ dùng để chỉ Quốc hội.}

\vocabentry{Partisan}
{/ˈpɑːrtɪzn/ (adj)}
{Prejudiced in favor of a particular cause or side.}
{The debate was marred by partisan politics rather than focusing on the facts.}
{Đây là một tính từ dùng để chỉ tính đảng phái hoặc sự ủng hộ mù quáng cho một phe.}

\vocabentry{Patriotism}
{/ˈpeɪtriətɪzəm/ (n)}
{The quality of being patriotic; devotion to and vigorous support for one's country.}
{Patriotism can be a powerful force for national unity during times of crisis.}
{Đây là một danh từ dùng để chỉ lòng yêu nước.}

\vocabentry{Policy}
{/ˈpɑːləsi/ (n)}
{A course or principle of action adopted or proposed by a government, party, business, or individual.}
{The government's foreign policy aims to promote peace and stability in the region.}
{Đây là một danh từ dùng để chỉ chính sách.}

\vocabentry{Politician}
{/ˌpɑːləˈtɪʃn/ (n)}
{A person who is professionally involved in politics, especially as a holder of an elected office.}
{Politicians need to listen to the concerns of their constituents to be successful.}
{Đây là một danh từ dùng để chỉ chính trị gia.}

\vocabentry{Polling}
{/ˈpoʊlɪŋ/ (n)}
{The recording of the votes of a body of people.}
{Recent polling suggests that the election will be very close.}
{Đây là một danh từ dùng để chỉ việc thăm dò ý kiến cử tri hoặc việc bỏ phiếu.}

\vocabentry{Protest}
{/ˈproʊtest/ (n)}
{A statement or action expressing disapproval of or objection to something.}
{Thousands of people took part in a peaceful protest against the new tax laws.}
{Đây là một danh từ dùng để chỉ sự phản kháng hoặc cuộc biểu tình.}

\vocabentry{Ratify}
{/ˈrætɪfaɪ/ (v)}
{Sign or give formal consent to (a treaty, contract, or agreement), making it officially valid. (Đây là một động từ dùng để chỉ việc thông qua hoặc phê chuẩn (hiệp ước, điều luật).).}
{The Senate is expected to ratify the international climate agreement next month.}
{}

\vocabentry{Referendum}
{/ˌrefəˈrendəm/ (n)}
{A general vote by the electorate on a single political question which has been referred to them for a direct decision.}
{The government has promised to hold a referendum on whether to change the voting system.}
{Đây là một danh từ dùng để chỉ cuộc trưng cầu dân ý.}

\vocabentry{Regime}
{/reɪˈʒiːm/ (n)}
{A government, especially an authoritarian one.}
{The old military regime was replaced by a democratically elected government.}
{Đây là một danh từ dùng để chỉ chế độ chính trị, thường mang nghĩa tiêu cực về sự độc tài.}

\vocabentry{Republic}
{/rɪˈpʌblɪk/ (n)}
{A state in which supreme power is held by the people and their elected representatives, and which has an elected president.}
{France is a republic with a strong central government and an elected president.}
{Đây là một danh từ dùng để chỉ nước cộng hòa.}

\vocabentry{Sanction}
{/ˈsæŋkʃn/ (n)}
{A threatened penalty for disobeying a law or rule.}
{The UN imposed economic sanctions on the country to pressure it to stop its nuclear program.}
{Đây là một danh từ dùng để chỉ sự trừng phạt hoặc lệnh cấm vận quốc tế.}

\vocabentry{Sovereignty}
{/ˈsɑːvrənti/ (n)}
{Supreme power or authority; the authority of a state to govern itself or another state.}
{The protection of national sovereignty is a primary goal of the defense policy.}
{Đây là một danh từ dùng để chỉ chủ quyền quốc gia.}

\vocabentry{Statute}
{/ˈstætʃuːt/ (n)}
{A written law passed by a legislative body.}
{The protection of human rights is enshrined in national statute.}
{Đây là một danh từ dùng để chỉ quy chế hoặc đạo luật ban hành bằng văn bản.}

\vocabentry{Suffrage}
{/ˈsʌfrɪdʒ/ (n)}
{The right to vote in political elections.}
{Universal suffrage was a major goal of the civil rights movement.}
{Đây là một danh từ dùng để chỉ quyền bầu cử.}

\vocabentry{Taxation}
{/tækˈseɪʃn/ (n)}
{The levying of tax.}
{Progressive taxation is often used to reduce income inequality in a society.}
{Đây là một danh từ dùng để chỉ hệ thống thuế hoặc việc đánh thuế.}

\vocabentry{Treaty}
{/ˈtriːti/ (n)}
{A formally concluded and ratified agreement between countries.}
{The two nations signed a peace treaty to end the long-running border dispute.}
{Đây là một danh từ dùng để chỉ hiệp ước quốc tế.}

\vocabentry{Tyranny}
{/ˈtɪrəni/ (n)}
{Cruel and oppressive government or rule.}
{The people rose up against the tyranny of the military dictatorship.}
{Đây là một danh từ dùng để chỉ sự bạo chính hoặc chế độ chuyên chế tàn bạo.}

\vocabentry{Unconstitutional}
{/ˌʌnkɑːnstɪˈtuːʃənl/ (adj)}
{Not in accordance with a political constitution.}
{The Supreme Court ruled that the new law was unconstitutional and must be struck down.}
{Đây là một tính từ dùng để chỉ những gì vi hiến hoặc trái với hiến pháp.}

\vocabentry{Veto}
{/ˈviːtoʊ/ (n/v)}
{A constitutional right to reject a decision or proposal made by a law-making body.}
{The President used his power of veto to block the controversial bill.}
{Đây là một danh từ/động từ dùng để chỉ quyền phủ quyết.}

\vocabentry{Welfare}
{/ˈwelfer/ (n)}
{Statutory procedure or social effort designed to promote the basic physical and material well-being of people in need.}
{The government provides welfare support for low-income families and the elderly.}
{Đây là một danh từ dùng để chỉ phúc lợi xã hội.}

\vocabentry{Appoint}
{/əˈpɔɪnt/ (v)}
{Assign a job or role to (someone).}
{The Prime Minister has the power to appoint senior judges and officials.}
{Đây là một động từ dùng để chỉ việc bổ nhiệm một chức vụ chính trị.}

\vocabentry{Civic}
{/ˈsɪvɪk/ (adj)}
{Relating to a city or town, especially its administration; municipal.}
{Civic engagement is essential for a vibrant and healthy community.}
{Đây là một tính từ dùng để chỉ nghĩa vụ công dân hoặc thuộc về đô thị.}

\vocabentry{Civil servant}
{/ˌsɪvl ˈsɜːrvənt/ (n)}
{A member of the civil service; a government official.}
{Civil servants play a vital role in implementing government policies and programs.}
{Đây là một danh từ dùng để chỉ công chức nhà nước.}

\vocabentry{Compulsory}
{/kəmˈpʌlsəri/ (adj)}
{Required by law or a rule; obligatory.}
{In some countries, voting in national elections is compulsory for all citizens.}
{Đây là một tính từ dùng để chỉ những gì bắt buộc theo luật pháp, ví dụ như bỏ phiếu bắt buộc.}

\vocabentry{Conservative}
{/kənˈsɜːrvətɪv/ (adj)}
{A person who is averse to change and holds to traditional values and attitudes.}
{The conservative party emphasizes the importance of family values and lower taxes.}
{Đây là một tính từ dùng để chỉ phe bảo thủ.}

\vocabentry{Diplomat}
{/ˈdɪpləmæt/ (n)}
{An official representing a country abroad.}
{A diplomat's job requires excellent communication and negotiation skills.}
{Đây là một danh từ dùng để chỉ nhà ngoại giao.}

\vocabentry{Dissolution}
{/ˌdɪsəˈluːʃn/ (n)}
{The closing down or dismissal of an assembly, partnership, or official body. (Đây là một danh từ dùng để chỉ sự giải tán (quốc hội, chính phủ).).}
{The dissolution of parliament signaled the start of a new election campaign.}
{}

\vocabentry{Federalism}
{/ˈfedərəlɪzəm/ (n)}
{The federal principle or system of government.}
{Federalism allows for a division of power between national and local governments.}
{Đây là một danh từ dùng để chỉ chủ nghĩa liên bang.}

\vocabentry{Impeachment}
{/ɪmˈpiːtʃmənt/ (n)}
{A charge of misconduct made against the holder of a public office.}
{The President faced impeachment proceedings following allegations of corruption.}
{Đây là một danh từ dùng để chỉ sự luận tội hoặc quy trình bãi miễn chức vụ.}

\vocabentry{International relations}
{/ˌɪntərˈnæʃnəl rɪˈleɪʃnz/ (n)}
{The academic study of interaction between nations.}
{She decided to major in international relations to pursue a career in diplomacy.}
{Đây là một danh từ dùng để chỉ ngành quan hệ quốc tế.}

\vocabentry{Jurisdiction}
{/ˌdʒʊrɪsˈdɪkʃn/ (n)}
{The official power to make legal decisions and judgments.}
{The crime falls under the jurisdiction of the federal authorities.}
{Đây là một danh từ dùng để chỉ thẩm quyền tài phán hoặc quyền hạn pháp lý.}

\vocabentry{Mayor}
{/ˈmeɪər/ (n)}
{The elected head of a city, town, or other municipality.}
{The mayor announced plans to build a new public library in the city center.}
{Đây là một danh từ dùng để chỉ thị trưởng.}

\vocabentry{Negotiate}
{/nəˈɡoʊʃieɪt/ (v)}
{Try to reach an agreement or compromise by discussion with others.}
{The two countries are trying to negotiate a new trade agreement.}
{Đây là một động từ dùng để chỉ việc đàm phán chính trị.}

\vocabentry{Nominate}
{/ˈnɑːmɪneɪt/ (v)}
{Propose or formally enter as a candidate for election or for an honor or award.}
{The party decided to nominate a fresh face for the upcoming election.}
{Đây là một động từ dùng để chỉ việc đề cử.}

\vocabentry{Official}
{/əˈfɪʃl/ (n/adj)}
{A person holding public office or having official duties.}
{Government officials are required to disclose their financial interests.}
{Đây là một danh từ/tính từ dùng để chỉ quan chức hoặc mang tính chính thức.}

\vocabentry{Petitioner}
{/pəˈtɪʃənər/ (n)}
{A person who presents a petition to an authority in respect of a particular cause.}
{The petitioners gathered thousands of signatures to oppose the new development.}
{Đây là một danh từ dùng để chỉ người kiến nghị.}

\vocabentry{Pluralism}
{/ˈplʊrəlɪzəm/ (n)}
{A condition or system in which two or more states, groups, principles, sources of authority, etc., coexist.}
{Political pluralism is essential for ensuring that different voices are heard in a democracy.}
{Đây là một danh từ dùng để chỉ sự đa nguyên chính trị.}

\vocabentry{Political party}
{/pəˈlɪtɪkl ˈpɑːrti/ (n)}
{An organized group of people who have the same political aims and opinions.}
{Joining a political party is a great way to get involved in local government.}
{Đây là một cụm danh từ dùng để chỉ đảng phái chính trị.}

\vocabentry{Presidency}
{/ˈprezɪdənsi/ (n)}
{The office of president.}
{His presidency was marked by significant economic growth and social change.}
{Đây là một danh từ dùng để chỉ chức danh tổng thống hoặc nhiệm kỳ tổng thống.}

\vocabentry{Radical}
{/ˈrædɪkl/ (adj)}
{Advocating or based on thorough or complete political or social reform.}
{The candidate's radical ideas attracted a lot of attention from the media.}
{Đây là một tính từ dùng để chỉ tư tưởng cấp tiến hoặc cực đoan.}

\vocabentry{Reform}
{/rɪˈfɔːrm/ (n/v)}
{Make changes in (something, typically a social, political, or economic institution or practice) in order to improve it.}
{The government is planning a major reform of the national healthcare system.}
{Đây là một danh từ/động từ dùng để chỉ sự cải cách.}

\vocabentry{Representative}
{/ˌreprɪˈzentətɪv/ (n)}
{A person chosen or appointed to act or speak for others.}
{Citizens elect representatives to parliament to voice their concerns.}
{Đây là một danh từ dùng để chỉ người đại diện hoặc hạ nghị sĩ.}

\vocabentry{Resignation}
{/ˌrezɪɡˈneɪʃn/ (n)}
{An act of retiring or giving up a position.}
{The minister's resignation came after weeks of intense pressure from the public.}
{Đây là một danh từ dùng để chỉ sự từ chức.}

\vocabentry{Resolution}
{/ˌrezəˈluːʃn/ (n)}
{A firm decision to do or not to do something.}
{The UN Security Council passed a resolution calling for an immediate ceasefire.}
{Trong chính trị là một nghị quyết.}

\vocabentry{Secular}
{/ˈsekjələr/ (adj)}
{Denoting attitudes, activities, or other things that have no religious or spiritual basis.}
{Turkey is a secular state where religion and politics are kept separate.}
{Đây là một tính từ dùng để chỉ tính thế tục, tách rời tôn giáo.}

\vocabentry{Senate}
{/ˈsenət/ (n)}
{The smaller upper assembly in the US Congress, most US state legislatures, and France.}
{The bill must be approved by the Senate before it can become law.}
{Đây là một danh từ dùng để chỉ Thượng viện.}

\vocabentry{Stability}
{/stəˈbɪləti/ (n)}
{The state of being stable.}
{Political stability is essential for attracting foreign investment.}
{Đây là một danh từ dùng để chỉ sự ổn định chính trị.}

\vocabentry{State}
{/steɪt/ (n)}
{A nation or territory considered as an organized political community under one government.}
{The state provides basic services like education, healthcare, and infrastructure.}
{Đây là một danh từ dùng để chỉ nhà nước hoặc bang.}

\vocabentry{Subsidize}
{/ˈsʌbsɪdaɪz/ (v)}
{Support (an organization or activity) financially.}
{The government subsidizes public transportation to make it more affordable for everyone.}
{Đây là một động từ dùng để chỉ việc trợ cấp từ ngân sách nhà nước.}

\vocabentry{Summit}
{/ˈsʌmɪt/ (n)}
{A meeting between heads of government.}
{The G7 summit focused on global economic issues and climate change.}
{Đây là một danh từ dùng để chỉ hội nghị thượng đỉnh.}

\vocabentry{Territory}
{/ˈterətɔːri/ (n)}
{An area of land under the jurisdiction of a ruler or state.}
{The dispute over the island territory has lasted for decades.}
{Đây là một danh từ dùng để chỉ lãnh thổ quốc gia.}

\vocabentry{Theocracy}
{/θiˈɑːkrəsi/ (n)}
{A system of government in which priests rule in the name of God or a god.}
{In a theocracy, religious laws often take precedence over civil laws.}
{Đây là một danh từ dùng để chỉ chính thể thần quyền.}

\vocabentry{Totalitarian}
{/toʊˌtæləˈteriən/ (adj)}
{Relating to a system of government that is centralized and dictatorial.}
{Life under a totalitarian regime is characterized by a complete lack of personal freedom.}
{Đây là một tính từ dùng để chỉ chế độ tổng trị, kiểm soát mọi mặt đời sống.}

\vocabentry{Turnout}
{/ˈtɜːrnaʊt/ (n)}
{The number of people who attend or participate in an event and especially the number of people who vote in an election.}
{The high voter turnout was seen as a sign of strong engagement with the political process.}
{Đây là một danh từ dùng để chỉ tỉ lệ cử tri đi bỏ phiếu.}

\vocabentry{Unanimous}
{/juˈnænɪməs/ (adj)}
{(Of two or more people) fully in agreement.}
{The committee's decision to approve the new project was unanimous.}
{Đây là một tính từ dùng để chỉ sự nhất trí hoàn toàn.}

\vocabentry{Underestimate}
{/ˌʌndərˈestɪmeɪt/ (v)}
{Estimate (something) to be smaller or less important than it actually is.}
{You should never underestimate the power of the people to bring about change.}
{Thường dùng cho việc đánh giá thấp sức mạnh chính trị hoặc đối thủ.}

\vocabentry{Universal}
{/ˌjuːnɪˈvɜːrsl/ (adj)}
{Of, affecting, or done by all people or things in the world or in a particular group.}
{The government is committed to providing universal access to quality education.}
{Ví dụ: universal healthcare - y tế toàn dân.}

\vocabentry{Uprising}
{/ˈʌpraɪzɪŋ/ (n)}
{An act of resistance or rebellion; a revolt.}
{The popular uprising eventually led to the downfall of the old regime.}
{Đây là một danh từ dùng để chỉ cuộc khởi nghĩa hoặc nổi dậy.}

\vocabentry{Urbanization}
{/ˌɜːrbənəˈzeɪʃn/ (n)}
{The process of making an area more urban.}
{Rapid urbanization presents many challenges for local governments in developing countries.}
{Ảnh hưởng đến quy hoạch chính trị và chính sách xã hội.}

\vocabentry{Utilization}
{/ˌjuːtələˈzeɪʃn/ (n)}
{The action of making practical and effective use of something.}
{We need to ensure the efficient utilization of public funds.}
{Trong quản lý tài chính công.}

\vocabentry{Value}
{/ˈvæljuː/ (n)}
{The regard that something is held to deserve.}
{Democratic values like freedom and equality are essential for a just society.}
{Giá trị chính trị cốt lõi.}

\vocabentry{Variable}
{/ˈveriəbl/ (n)}
{An element, feature, or factor that is liable to vary or change.}
{Economic performance is a key variable in determining the outcome of an election.}
{Các biến số trong phân tích chính trị.}

\vocabentry{Vibrant}
{/ˈvaɪbrənt/ (adj)}
{Full of energy and life.}
{A vibrant civil society is essential for holding the government accountable.}
{Một xã hội chính trị sôi động.}

\vocabentry{Virtual}
{/ˈvɜːrtʃuəl/ (adj)}
{Almost or nearly as described, but not completely.}
{The candidate held a virtual town hall meeting to reach voters during the pandemic.}
{Trong chính trường số.}

\vocabentry{Visibility}
{/ˌvɪzəˈbɪləti/ (n)}
{The state of being able to see or be seen.}
{High visibility is essential for politicians who want to build a strong public profile.}
{Sự công khai, minh bạch của quan chức.}

\vocabentry{Vital}
{/ˈvaɪtl/ (adj)}
{Absolutely necessary or important; essential.}
{National unity is vital for overcoming the challenges facing the country.}
{Cần thiết cho vận mệnh quốc gia.}

\vocabentry{Volatility}
{/ˌvɑːləˈtɪləti/ (n)}
{Liability to change rapidly and unpredictably, especially for the worse.}
{Political volatility can discourage investors and harm the economy.}
{Sự biến động của thị trường chính trị.}

\vocabentry{Volume}
{/ˈvɑːljuːm/ (n)}
{The amount of space that a substance or object occupies.}
{The sheer volume of votes to be counted meant that the result was delayed.}
{Trong thống kê bầu cử.}

\vocabentry{Voluntary}
{/ˈvɑːlənteri/ (adj)}
{Done, given, or acting of one's own free will.}
{Many citizens engage in voluntary work to support their local community.}
{Sự tham gia chính trị tự nguyện.}

\vocabentry{Vulnerability}
{/ˌvʌlnərəˈbɪləti/ (n)}
{The quality or state of being exposed to the possibility of being attacked or harmed.}
{The country's vulnerability to cyberattacks is a major concern for national security.}
{Sự dễ bị tổn thương của nền an ninh.}

\vocabentry{Welfare state}
{/ˈwelfer steɪt/ (n)}
{A system whereby the state undertakes to protect the health and well-being of its citizens.}
{The creation of the welfare state helped to reduce poverty and improve living standards.}
{Đây là một danh từ dùng để chỉ nhà nước phúc lợi.}

\vocabentry{Wholesale}
{/ˈhoʊlseɪl/ (adj)}
{Done on a large scale; extensive.}
{The government's wholesale reform of the tax system was met with mixed reactions.}
{Sự thay đổi chính sách quy mô lớn.}

\vocabentry{Widespread}
{/ˌwaɪdˈspred/ (adj)}
{Found or distributed over a large area or number of people.}
{There is widespread support for the new environmental regulations.}
{Sự ủng hộ rộng rãi.}

\vocabentry{111. Abuse of power}
{/əˈbjuːs əv ˈpaʊər/ (n)}
{The misuse of a position of power to take advantage of others or to achieve personal gain.}
{The investigation revealed a widespread abuse of power within the local police department.}
{Đây là một cụm danh từ dùng để chỉ hành vi lạm dụng quyền lực hoặc chức vụ để trục lợi cá nhân.}

\vocabentry{112. Activism}
{/ˈæktɪvɪzəm/ (n)}
{The policy or action of using vigorous campaigning to bring about political or social change.}
{Student activism played a crucial role in the movement for environmental protection laws.}
{Đây là một danh từ dùng để chỉ chủ nghĩa hoạt động, nỗ lực thúc đẩy thay đổi chính trị xã hội.}

\vocabentry{113. Agenda}
{/əˈdʒendə/ (n)}
{A list of items to be discussed at a formal meeting; also a set of goals or plans for a political group.}
{Education reform is at the top of the government's political agenda for the next year.}
{Đây là một danh từ dùng để chỉ chương trình nghị sự hoặc kế hoạch hành động chính trị.}

\vocabentry{114. Alliance}
{/əˈlaɪəns/ (n)}
{A union or association formed for mutual benefit, especially between countries or organizations.}
{The two countries formed a strategic alliance to strengthen their defense against potential threats.}
{Đây là một danh từ dùng để chỉ sự liên minh hoặc đồng minh giữa các quốc gia.}

\vocabentry{115. Appropriation}
{/əˌproʊpriˈeɪʃn/ (n)}
{The action of devoting money to a particular purpose, especially by a government.}
{The senate approved an extra appropriation of funds for the national healthcare system.}
{Đây là một danh từ dùng để chỉ sự cấp ngân sách hoặc phân bổ tài chính cho một mục đích cụ thể.}

\vocabentry{116. Arbitrary}
{/ˈɑːrbɪtreri/ (adj)}
{Based on random choice or personal whim, rather than any reason or system.}
{The protesters criticized the government for making arbitrary arrests of opposition leaders.}
{Đây là một tính từ dùng để chỉ sự độc đoán hoặc tùy tiện trong việc ra quyết định chính trị.}

\vocabentry{117. Armistice}
{/ˈɑːrmɪstɪs/ (n)}
{An agreement made by opposing sides in a war to stop fighting for a certain time; a truce.}
{After years of conflict, both nations finally signed an armistice to discuss peace terms.}
{Đây là một danh từ dùng để chỉ sự đình chiến hoặc thỏa thuận ngừng bắn.}

\vocabentry{118. Asylum}
{/əˈsaɪləm/ (n)}
{Protection granted by a nation to someone who has left their native country as a political refugee.}
{The journalist sought political asylum in Europe after being threatened in his home country.}
{Đây là một danh từ dùng để chỉ sự tị nạn chính trị.}

\vocabentry{119. Authoritarianism}
{/əˌθɔːrəˈteriənɪzəm/ (n)}
{The enforcement or advocacy of strict obedience to authority at the expense of personal freedom.}
{The rise of authoritarianism in the region has led to a significant crackdown on free speech.}
{Đây là một danh từ dùng để chỉ chủ nghĩa uy quyền hoặc chế độ độc tài.}

\vocabentry{120. Autonomy}
{/ɔːˈtɑːnəmi/ (n)}
{The right or condition of self-government.}
{The province is seeking greater autonomy from the central government in managing its resources.}
{Đây là một danh từ dùng để chỉ quyền tự trị của một khu vực hoặc tổ chức.}

\vocabentry{121. Backbencher}
{/ˌbækˈbentʃər/ (n)}
{A member of parliament who does not hold office in the government or opposition and sits on the back benches.}
{Several backbenchers spoke out against the proposed changes to the immigration law.}
{Đây là một danh từ dùng để chỉ những nghị sĩ không giữ chức vụ then chốt trong chính phủ.}

\vocabentry{122. Ballot box}
{/ˈbælət bɑːks/ (n)}
{A sealed box into which voters put completed ballot papers.}
{The security of the ballot boxes is essential for ensuring a fair and transparent election.}
{Đây là một cụm danh từ dùng để chỉ hòm phiếu bầu.}

\vocabentry{123. Biased}
{/ˈbaɪəst/ (adj)}
{Unfairly prejudiced for or against someone or something.}
{The report was criticized for being heavily biased in favor of the current administration.}
{Đây là một tính từ dùng để chỉ sự thiên vị, thiếu khách quan trong truyền thông hoặc chính trị.}

\vocabentry{124. Border control}
{/ˈbɔːrdər kənˈtroʊl/ (n)}
{The measures taken by a country to monitor and regulate its borders.}
{The government has announced stricter border control measures to curb illegal immigration.}
{Đây là một cụm danh từ dùng để chỉ hệ thống kiểm soát biên giới quốc gia.}

\vocabentry{125. Bribery}
{/ˈbraɪbəri/ (n)}
{The giving or offering of a bribe.}
{Several officials were arrested on charges of bribery and money laundering.}
{Đây là một danh từ dùng để chỉ hành vi hối lộ, một vấn nạn trong chính trường.}

\vocabentry{126. Budget deficit}
{/ˈbʌdʒɪt ˈdefɪsɪt/ (n)}
{An indicator of financial health in which expenditures exceed revenue.}
{Reducing the budget deficit is a top priority for the finance minister this year.}
{Đây là một cụm danh từ dùng để chỉ tình trạng thâm hụt ngân sách chính phủ.}

\vocabentry{127. By-election}
{/ˈbaɪ ɪˌlekʃn/ (n)}
{An election held in a single political constituency to fill a vacancy.}
{The opposition party won a surprising victory in the recent by-election.}
{Đây là một danh từ dùng để chỉ cuộc bầu cử bổ sung khi có một ghế nghị sĩ bị bỏ trống.}

\vocabentry{128. Canvassing}
{/ˈkænvəsɪŋ/ (n)}
{The systematic initiation of direct contact with individuals, commonly used during political campaigns.}
{Volunteers spent the weekend canvassing the neighborhood to gain support for their candidate.}
{Đây là một danh từ dùng để chỉ hoạt động vận động tranh cử tận nhà cử tri.}

\vocabentry{129. Capitalism}
{/ˈkæpɪtəlɪzəm/ (n)}
{An economic and political system in which a country's trade and industry are controlled by private owners.}
{Some critics argue that unchecked capitalism can lead to extreme wealth inequality.}
{Đây là một danh từ dùng để chỉ chủ nghĩa tư bản.}

\vocabentry{130. Censure}
{/ˈsenʃər/ (n)}
{The expression of formal disapproval.}
{The minister faced a vote of censure following the failure of the new policy.}
{Đây là một danh từ dùng để chỉ sự khiển trách hoặc chỉ trích chính thức trong quốc hội hoặc tổ chức.}

\vocabentry{131. Centralization}
{/ˌsentrələˈzeɪʃn/ (n)}
{The concentration of control of an activity or organization under a single authority.}
{The centralization of power has led to concerns about the loss of local decision-making.}
{Đây là một danh từ dùng để chỉ sự tập quyền, dồn quyền lực về trung ương.}

\vocabentry{132. Checks and balances}
{/ˌtʃeks ənd ˈbælənsɪz/ (n)}
{Counterbalancing influences by which an organization or system is regulated.}
{The system of checks and balances ensures that no single branch of government becomes too powerful.}
{Đây là một cụm danh từ dùng để chỉ hệ thống kiểm soát và đối trọng quyền lực.}

\vocabentry{133. Citizen}
{/ˈsɪtɪzn/ (n)}
{A legally recognized subject or national of a state or commonwealth.}
{Every citizen has the right to vote and the duty to obey the law.}
{Đây là một danh từ dùng để chỉ công dân của một quốc gia.}

\vocabentry{134. Civil liberties}
{/ˌsɪvl ˈlɪbərtiz/ (n)}
{Being subject only to laws established for the good of the community, especially with regard to freedom of action and speech.}
{Civil liberties must be protected even during times of national emergency.}
{Đây là một cụm danh từ dùng để chỉ các quyền tự do dân sự.}

\vocabentry{135. Civil rights}
{/ˌsɪvl ˈraɪts/ (n)}
{The rights of citizens to political and social freedom and equality.}
{The civil rights movement fought to end racial discrimination and achieve equality.}
{Đây là một cụm danh từ dùng để chỉ các quyền dân sự.}

\vocabentry{136. Civil war}
{/ˌsɪvl ˈwɔːr/ (n)}
{A war between citizens of the same country.}
{The country was devastated by a long and brutal civil war that lasted over a decade.}
{Đây là một cụm danh từ dùng để chỉ nội chiến.}

\vocabentry{137. Cold war}
{/ˌkoʊld ˈwɔːr/ (n)}
{A state of political hostility between countries characterized by threats and propaganda.}
{The Cold War era was marked by intense competition between the US and the Soviet Union.}
{Đây là một cụm danh từ dùng để chỉ chiến tranh lạnh.}

\vocabentry{138. Colonialism}
{/kəˈloʊniəlɪzəm/ (n)}
{The policy or practice of acquiring full or partial political control over another country.}
{The effects of colonialism are still felt in many former colonies around the world.}
{Đây là một danh từ dùng để chỉ chủ nghĩa thực dân.}

\vocabentry{139. Communism}
{/ˈkɑːmjunɪzəm/ (n)}
{A political theory derived from Karl Marx, advocating class war and leading to a society in which all property is publicly owned.}
{The fall of the Berlin Wall symbolized the decline of communism in Eastern Europe.}
{Đây là một danh từ dùng để chỉ chủ nghĩa cộng sản.}

\vocabentry{140. Compliance}
{/kəmˈplaɪəns/ (n)}
{The action or fact of complying with a wish or command.}
{The government is monitoring industrial compliance with the new environmental laws.}
{Trong chính trị, đây là sự tuân thủ các quy định pháp luật.}

\vocabentry{141. Concession}
{/kənˈseʃn/ (n)}
{A thing that is granted, especially in response to demands.}
{Both sides had to make significant concessions to reach a peaceful resolution.}
{Đây là một danh từ dùng để chỉ sự nhượng bộ chính trị để đạt được thỏa thuận.}

\vocabentry{142. Conflict of interest}
{/ˌkɑːnflɪkt əv ˈɪntrest/ (n)}
{A situation in which a person is in a position to derive personal benefit from actions made in their official capacity.}
{The official was accused of a conflict of interest for awarding the contract to his own firm.}
{Đây là một cụm danh từ dùng để chỉ sự xung đột lợi ích.}

\vocabentry{143. Congress}
{/ˈkɑːŋɡres/ (n)}
{A formal meeting or series of meetings for discussion between delegates.}
{The new bill is currently being debated by members of Congress.}
{Đây là một danh từ dùng để chỉ Quốc hội, đặc biệt là ở Hoa Kỳ.}

\vocabentry{144. Consensus}
{/kənˈsensəs/ (n)}
{A general agreement.}
{There is a growing consensus among lawmakers that the tax system needs to be reformed.}
{Đây là một danh từ dùng để chỉ sự đồng thuận chung giữa các bên chính trị.}

\vocabentry{145. Coup d'état}
{/ˌkuː deɪˈtɑː/ (n)}
{A sudden, violent, and illegal seizure of power from a government.}
{The military leader led a successful coup d'état and took control of the country.}
{Đây là một danh từ dùng để chỉ cuộc đảo chính quân sự hoặc lật đổ chính quyền.}

\vocabentry{146. Delegation}
{/ˌdelɪˈɡeɪʃn/ (n)}
{A body of delegates or representatives.}
{A delegation of senior officials was sent to negotiate the peace treaty.}
{Đây là một danh từ dùng để chỉ phái đoàn đại biểu tham dự hội nghị.}

\vocabentry{147. Demagogue}
{/ˈdeməɡɑːɡ/ (n)}
{A political leader who seeks support by appealing to the desires and prejudices of ordinary people.}
{Critics accused the candidate of being a demagogue who used fear to gain votes.}
{Đây là một danh từ dùng để chỉ kẻ mị dân, lôi kéo cử tri bằng cảm xúc hơn là lý trí.}

\vocabentry{148. Depoliticize}
{/ˌdiːpəˈlɪtɪsaɪz/ (v)}
{Remove from political influence or focus.}
{There are calls to depoliticize the process of appointing senior judges.}
{Đây là một động từ dùng để chỉ việc phi chính trị hóa một vấn đề nào đó.}

\vocabentry{149. Diplomatic immunity}
{/ˌdɪpləmætɪk ɪˈmjuːnəti/ (n)}
{The privilege of exemption from certain laws and taxes granted to diplomats.}
{Diplomatic immunity protects foreign officials from being prosecuted in the host country.}
{Đây là một cụm danh từ dùng để chỉ quyền miễn trừ ngoại giao.}

\vocabentry{150. Disenfranchisement}
{/ˌdɪsɪnˈfræntʃaɪzmənt/ (n)}
{The state of being deprived of a right or privilege, especially the right to vote.}
{The new voting regulations have led to the disenfranchisement of many minority voters.}
{Đây là một danh từ dùng để chỉ sự tước quyền công dân hoặc quyền bầu cử.}

\vocabentry{151. Dissent}
{/dɪˈsent/ (n)}
{The expression or holding of opinions at variance with those previously or commonly or officially held.}
{The government has been accused of attempting to suppress all forms of political dissent.}
{Đây là một danh từ dùng để chỉ sự bất đồng quan điểm chính trị.}

\vocabentry{152. Domestic affairs}
{/dəˌmestɪk əˈferz/ (n)}
{Issues or concerns that are internal to a country.}
{The President spent much of the speech focusing on domestic affairs like education and jobs.}
{Đây là một cụm danh từ dùng để chỉ các vấn đề đối nội của một quốc gia.}

\vocabentry{153. Dominion}
{/dəˈmɪnjən/ (n)}
{Sovereignty or control.}
{The historical records detail the extent of the empire's colonial dominion.}
{Đây là một danh từ dùng để chỉ quyền thống trị hoặc lãnh thổ thuộc quyền kiểm soát.}

\vocabentry{154. Egalitarianism}
{/iˌɡælɪˈteriənɪzəm/ (n)}
{The doctrine that all people are equal and deserve equal rights and opportunities.}
{His political philosophy is based on the principles of egalitarianism and social justice.}
{Đây là một danh từ dùng để chỉ chủ nghĩa bình quân.}

\vocabentry{155. Embargo}
{/ɪmˈbɑːrɡoʊ/ (n)}
{An official ban on trade or other commercial activity with a particular country.}
{The international community imposed an arms embargo on the nation to stop the violence.}
{Đây là một danh từ dùng để chỉ lệnh cấm vận thương mại.}

\vocabentry{156. Enactment}
{/ɪˈnæktmənt/ (n)}
{The process of passing legislation.}
{The enactment of the new healthcare bill will provide coverage for millions of people.}
{Đây là một danh từ dùng để chỉ việc ban hành luật hoặc quá trình một đạo luật có hiệu lực.}

\vocabentry{157. Espionage}
{/ˈespɪənɑːʒ/ (n)}
{The practice of spying or of using spies, typically by governments to obtain political or military information.}
{Several individuals were charged with espionage for leaking military secrets.}
{Đây là một danh từ dùng để chỉ hoạt động gián điệp.}

\vocabentry{158. Establishment}
{/ɪˈstæblɪʃmənt/ (n)}
{The ruling class of a society; the authorities in power.}
{The populist candidate campaigned against the political establishment in the capital.}
{Trong chính trị, đây là danh từ chỉ tầng lớp cầm quyền hoặc giới tinh hoa chính trị.}

\vocabentry{159. Exile}
{/ˈeksaɪl/ (n)}
{The state of being barred from one's native country, typically for political or punitive reasons.}
{The former leader lived in exile for twenty years before returning to his country.}
{Đây là một danh từ dùng để chỉ sự lưu đày hoặc cuộc sống lưu vong.}

\vocabentry{160. Faction}
{/ˈfækʃn/ (n)}
{A small organized dissenting group within a larger one, especially in politics.}
{Rival factions within the party are struggling for control of the leadership.}
{Đây là một danh từ dùng để chỉ phe phái trong một tổ chức hoặc đảng phái.}

\vocabentry{161. Foreign policy}
{/ˈfɔːrən ˈpɑːləsi/ (n)}
{A government's strategy in dealing with other nations.}
{The nation's foreign policy is focused on promoting trade and regional security.}
{Đây là một cụm danh từ dùng để chỉ chính sách đối ngoại.}

\vocabentry{162. Frontier}
{/frʌnˈtɪr/ (n)}
{A line or border separating two countries.}
{Thousands of refugees are gathered at the frontier, waiting to enter the neighboring country.}
{Đây là một danh từ dùng để chỉ biên giới quốc gia.}

\vocabentry{163. Gerrymandering}
{/ˈdʒerimændərɪŋ/ (n)}
{Manipulate the boundaries of (an electoral constituency) so as to favor one party or class.}
{Gerrymandering is often used to create "safe seats" for a particular political party.}
{Đây là một danh từ dùng để chỉ việc gian lận trong phân chia khu vực bầu cử.}

\vocabentry{164. Globalism}
{/ˈɡloʊbəlɪzəm/ (n)}
{The operation or planning of economic and foreign policy on a global basis.}
{Advocates of globalism argue that international cooperation is essential for solving world problems.}
{Đây là một danh từ dùng để chỉ chủ nghĩa toàn cầu.}

\vocabentry{165. Grassroots}
{/ˈɡræsruːts/ (n)}
{The most basic level of an activity or organization; ordinary people in a society or political party.}
{The campaign relied on grassroots support from thousands of local volunteers.}
{Đây là một danh từ dùng để chỉ tầng lớp bình dân hoặc phong trào từ cơ sở.}

\vocabentry{166. Hegemony}
{/hɪˈdʒeməni/ (n)}
{Leadership or dominance, especially by one country or social group over others.}
{The country's economic and military hegemony in the region is being challenged.}
{Đây là một danh từ dùng để chỉ quyền bá chủ hoặc sự thống trị.}

\vocabentry{167. Immigration policy}
{/ˌɪmɪˈɡreɪʃn ˈpɑːləsi/ (n)}
{The rules and regulations a government uses to manage immigration.}
{The government is revising its immigration policy to attract more skilled workers.}
{Đây là một cụm danh từ dùng để chỉ chính sách nhập cư.}

\vocabentry{168. Imperialism}
{/ɪmˈpɪriəlɪzəm/ (n)}
{A policy of extending a country's power and influence through diplomacy or military force.}
{The history of the 19th century was marked by European imperialism in Africa and Asia.}
{Đây là một danh từ dùng để chỉ chủ nghĩa đế quốc.}

\vocabentry{169. Industrialization}
{/ɪnˌdʌstriələˈzeɪʃn/ (n)}
{The development of industries in a country or region on a wide scale.}
{Rapid industrialization led to major shifts in political power and labor laws.}
{Đây là một danh từ ảnh hưởng mạnh đến cấu trúc chính trị và xã hội.}

\vocabentry{170. Infrastructure}
{/ˈɪnfrəstrʌktʃər/ (n)}
{The basic physical and organizational structures needed for the operation of a society.}
{The government is investing heavily in transport infrastructure to boost the economy.}
{Đây là một danh từ chỉ cơ sở hạ tầng mà chính phủ có trách nhiệm xây dựng.}

\vocabentry{171. Insurgency}
{/ɪnˈsɜːrdʒənsi/ (n)}
{An active revolt or uprising.}
{The military is struggling to contain the insurgency in the northern provinces.}
{Đây là một danh từ dùng để chỉ sự nổi dậy hoặc cuộc nổi loạn chống chính quyền.}

\vocabentry{172. Integration}
{/ˌɪntɪˈɡreɪʃn/ (n)}
{The action or process of integrating.}
{European integration has led to greater cooperation between member nations.}
{Trong chính trị là sự hội nhập quốc tế hoặc hòa nhập sắc tộc.}

\vocabentry{173. Intervention}
{/ˌɪntərˈvenʃn/ (n)}
{The action or process of intervening, especially in another state's affairs.}
{Many people opposed the military intervention in the neighboring country's civil war.}
{Đây là một danh từ dùng để chỉ sự can thiệp quốc tế.}

\vocabentry{174. Isolationism}
{/ˌaɪsəˈleɪʃənɪzəm/ (n)}
{A policy of remaining apart from the affairs or interests of other groups, especially the political affairs of other countries.}
{The candidate's isolationism appealed to voters who were tired of foreign wars.}
{Đây là một danh từ dùng để chỉ chủ nghĩa biệt lập.}

\vocabentry{175. Legislative}
{/ˈledʒɪsleɪtɪv/ (adj)}
{Having the power to make laws.}
{The legislative branch of the government is responsible for drafting new bills.}
{Đây là một tính từ dùng để chỉ nhánh hoặc quyền lập pháp.}

\vocabentry{176. Local government}
{/ˌloʊkl ˈɡʌvərnmənt/ (n)}
{The administration of a particular town, county, or district.}
{Local government is responsible for services like waste collection and libraries.}
{Đây là một cụm danh từ dùng để chỉ chính quyền địa phương.}

\vocabentry{177. Majority}
{/məˈdʒɔːrəti/ (n)}
{The greater number.}
{The party won a large majority in the national assembly, giving them total control.}
{Đây là một danh từ dùng để chỉ đa số phiếu bầu hoặc ghế trong quốc hội.}

\vocabentry{178. Mandate}
{/ˈmændeɪt/ (n)}
{An official order or commission to do something.}
{The election victory gave the new government a clear mandate for reform.}
{Trong chính trị là sự ủy nhiệm quyền lực từ cử tri.}

\vocabentry{179. Manifesto}
{/ˌmænɪˈfestoʊ/ (n)}
{A public declaration of policy and aims, especially one issued before an election by a political party or candidate.}
{The party's election manifesto included promises of lower taxes and better schools.}
{Đây là một danh từ dùng để chỉ bản tuyên ngôn của một đảng.}

\vocabentry{180. Martial law}
{/ˌmɑːrʃl ˈlɔː/ (n)}
{Military government involving the suspension of ordinary law.}
{The government declared martial law in response to the widespread violence in the capital.}
{Đây là một cụm danh từ dùng để chỉ thiết quân luật.}

\vocabentry{181. Mass media}
{/ˌmæs ˈmiːdiə/ (n)}
{The various methods of communication that reach large numbers of people.}
{The mass media plays a crucial role in shaping the political views of the public.}
{Trong chính trị là công cụ tuyên truyền và thông tin.}

\vocabentry{182. Minority government}
{/maɪˈnɔːrəti ˈɡʌvərnmənt/ (n)}
{A government in which the governing party has most seats but still less than half the total.}
{A minority government often struggles to pass new legislation without opposition support.}
{Đây là một cụm danh từ dùng để chỉ chính phủ thiểu số.}

\vocabentry{183. Modernization}
{/ˌmɑːrdənəˈzeɪʃn/ (n)}
{The process of adapting something to modern needs or habits.}
{The government is committed to the modernization of the national economy.}
{Chính sách hiện đại hóa đất nước.}

\vocabentry{184. Municipal}
{/mjuːˈnɪsɪpl/ (adj)}
{Relating to a city or town or its governing body.}
{Municipal elections will be held next month to choose a new city council.}
{Đây là một tính từ dùng để chỉ các cấp hành chính thuộc thành phố/đô thị.}

\vocabentry{185. Nationalism}
{/ˈnæʃnəlɪzəm/ (n)}
{Patriotic feeling, principles, or efforts.}
{The rise of aggressive nationalism has led to increased tensions between neighboring countries.}
{Đây là một danh từ dùng để chỉ chủ nghĩa dân tộc.}

\vocabentry{186. Neutrality}
{/nuːˈtræləti/ (n)}
{The state of not supporting or helping either side in a conflict, disagreement, etc.}
{The country has a long history of political and military neutrality in international conflicts.}
{Đây là một danh từ dùng để chỉ sự trung lập.}

\vocabentry{187. Oligarchy}
{/ˈɑːlɪɡɑːrki/ (n)}
{A small group of people having control of a country, organization, or institution.}
{Some critics argue that the country is ruled by a wealthy and powerful oligarchy.}
{Đây là một danh từ dùng để chỉ chính thể tập quyền trong tay một nhóm nhỏ.}

\vocabentry{188. Oppression}
{/əˈpreʃn/ (n)}
{Prolonged cruel or unjust treatment or control.}
{The struggle for freedom was a fight against decades of systematic political oppression.}
{Đây là một danh từ dùng để chỉ sự áp bức, kìm kẹp chính trị.}

\vocabentry{189. Overthrow}
{/ˌoʊvərˈθroʊ/ (v)}
{Remove forcibly from power.}
{There were rumors of a plot to overthrow the government and install a new leader.}
{Đây là một động từ dùng để chỉ việc lật đổ chính quyền.}

\vocabentry{190. Pacifism}
{/ˈpæsɪfɪzəm/ (n)}
{The belief that any violence, including war, is unjustifiable and that all disputes should be settled by peaceful means.}
{Her political views are strongly influenced by her commitment to pacifism.}
{Đây là một danh từ dùng để chỉ chủ nghĩa hòa bình.}

\vocabentry{191. Propaganda}
{/ˌprɑːpəˈɡændə/ (n)}
{Information, especially of a biased or misleading nature, used to promote a political cause or point of view.}
{The state-run media was accused of spreading government propaganda.}
{Lặp lại để nhấn mạnh.}

\vocabentry{192. Protocol}
{/ˈproʊtəkɑːl/ (n)}
{The official procedure or system of rules governing affairs of state or diplomatic occasions.}
{It is important to follow diplomatic protocol when hosting foreign heads of state.}
{Đây là một danh từ dùng để chỉ nghi thức ngoại giao.}

\vocabentry{193. Public sector}
{/ˌpʌblɪk ˈsektər/ (n)}
{The part of an economy that is controlled by the state.}
{Many people work in the public sector as teachers, police officers, and civil servants.}
{Đây là một cụm danh từ dùng để chỉ khu vực công.}

\vocabentry{194. Quorum}
{/ˈkwɔːrəm/ (n)}
{The minimum number of members of an assembly or society that must be present at any of its meetings to make the proceedings of that meeting valid.}
{The meeting had to be postponed because there was no quorum of members present.}
{Đây là một danh từ dùng để chỉ túc số - số đại biểu tối thiểu.}

\vocabentry{195. Rapprochement}
{/ˌræproʊʃˈmɑ̃ː/ (n)}
{(Especially in international relations) an establishment or resumption of harmonious relations.}
{The meeting between the two leaders led to a historic rapprochement between the nations.}
{Đây là một danh từ dùng để chỉ sự nối lại quan hệ hữu nghị.}

\vocabentry{196. Red tape}
{/ˌred ˈteɪp/ (n)}
{Excessive bureaucracy or adherence to rules and formalities, especially in public business.}
{Small businesses often struggle with the amount of government red tape they have to deal with.}
{Đây là một danh từ dùng để chỉ các thủ tục hành chính rườm rà.}

\vocabentry{197. Security Council}
{/sɪˈkjʊrəti ˈkaʊnsl/ (n)}
{A permanent council of the United Nations.}
{The Security Council is responsible for maintaining international peace and security.}
{Đây là một cụm danh từ dùng để chỉ Hội đồng Bảo an Liên Hợp Quốc.}

\vocabentry{198. Separation of powers}
{/ˌsepəˈreɪʃn əv ˈpaʊərz/ (n)}
{An act of vesting the legislative, executive, and judicial powers of government in separate bodies.}
{The separation of powers is a fundamental principle of the US government.}
{Đây là một cụm danh từ dùng để chỉ thuyết tam quyền phân lập.}

\vocabentry{199. Socialism}
{/ˈsoʊʃəlɪzəm/ (n)}
{A political and economic theory of social organization which advocates that the means of production, distribution, and exchange should be owned or regulated by the community as a whole.}
{Socialism emphasizes the importance of social equality and public services.}
{Đây là một danh từ dùng để chỉ chủ nghĩa xã hội.}

\vocabentry{200. Theocratic}
{/ˌθiːəˈkrætɪk/ (adj)}
{Relating to or denoting a system of government in which priests rule in the name of God or a god.}
{}
{Đây là một tính từ dùng để chỉ thuộc về chính thể thần quyền.}

